\chapter{Metryka i realizacja projektu}
\label{cha:LabSI}
\makeatletter

% *********************************************************************
% W poniższej metryczce należy uzupełnić informacje o:
% - temacie ćwiczenia
% - numerze ćwiczenia
% - numerze grupy laboratoryjnej
% - numerze zespołu
% - dacie wykonania oraz oddania ćwiczenia
% *********************************************************************

\begin{table}[H]
    \centering
    \renewcommand{\tabularxcolumn}[1]{m{#1}}
    \newcolumntype{C}{>{\centering\arraybackslash}X}
    \begin{tabularx}{\linewidth}{|C|C|}
        \hline
        \multicolumn{2}{|c|}{\makecell{\textbf{\Large{Sztuczna Inteligencja}} \\ \textbf{Sprawozdanie projektowe}}} \\ \hline
        \multicolumn{1}{|l|}{Temat} & Klasyfikacja znaków dłoni (Sign Language MNIST) \\ \hline
        \multicolumn{1}{|l|}{Ćwiczenie nr:} & Projekt semestralny \\ \hline
        \multicolumn{1}{|l|}{Autorzy:} & \@author \\ \hline
        \multicolumn{1}{|l|}{Grupa laboratoryjna} & XX (do uzupełnienia) \\ \hline
        \multicolumn{1}{|l|}{Zespół nr:} & XX (do uzupełnienia) \\ \hline
        \multicolumn{1}{|l|}{Data wykonania ćwiczenia} & czerwiec 2026 \\ \hline
        \multicolumn{1}{|l|}{Data oddania ćwiczenia} & czerwiec 2026 \\ \hline
    \end{tabularx}
\end{table}

\section{Opis realizacji}
W projekcie zaimplementowano aplikację CLI do klasyfikacji znaków dłoni na bazie
zbioru Sign Language MNIST. Aplikacja obsługuje trzy tryby pracy:
\texttt{train}, \texttt{eval} oraz \texttt{infer}. Architektura kodu została
podzielona na moduły danych, modelowania i inferencji.

\section{Wykorzystane narzędzia}
W realizacji wykorzystano:
\begin{itemize}
    \item Python oraz menedżer środowiska \texttt{uv},
    \item \texttt{numpy}, \texttt{pandas}, \texttt{scikit-learn},
    \item \texttt{tensorflow} dla modelu CNN,
    \item \texttt{opencv-python} i \texttt{mediapipe} dla inferencji kamerowej,
    \item \texttt{pytest} i \texttt{ruff} do walidacji jakości kodu.
\end{itemize}

\section{Zakres funkcjonalny}
Zaimplementowane funkcje systemu:
\begin{enumerate}
    \item trening modelu i zapis artefaktu,
    \item ewaluacja jakości na zbiorze testowym,
    \item inferencja z pojedynczego rekordu CSV,
    \item inferencja online z kamery z podglądem detekcji.
\end{enumerate}

\section{Wynik końcowy}
Otrzymano działający prototyp klasyfikatora znaków dłoni, który może być
uruchamiany lokalnie i rozszerzany o kolejne metody modelowania.
Dokumentację przygotowano w oparciu o rzeczywistą strukturę kodu i aktualny
stan repozytorium.
